% ------------------------------------------------------------------------
%
\begin{frame}{Tail call optimisation}

  Let us define a function that sums integers in a non\hyp{}empty
  list:
  \begin{equation*}
    \fun{sum} \, [n] \xrightarrow{\alpha} n\qquad
    \fun{sum}(\cons{n}{s}) \xrightarrow{\beta} n + \fun{sum}(s)
  \end{equation*}
  where we introduced a notation `\(\fun{sum}[n]\)' instead of
  `\(\fun{sum}([n])\)'. We will omit the parentheses when there is
  exactly one argument and it is clear when it starts and ends.

  \bigskip

  The view of the system with maximum sharing is:
  \begin{center}
    \includegraphics{sum_dag}
  \end{center}

\end{frame}

% ------------------------------------------------------------------------
%
\begin{frame}{Tail call optimisation}

  The first phase of the evaluation of \(\fun{sum}([3;7;5])\) is:
  \begin{center}
    \includegraphics[bb=72 635 404 721,scale=0.99]{sum375_push}
  \end{center}

\end{frame}

% ------------------------------------------------------------------------
%
\begin{frame}{Tail call optimisation}

  The second phase (replacing the pending calls with their values) is
  \begin{center}
    \includegraphics[bb=72 644 380 721]{sum375_pop}
  \end{center}

\end{frame}

% ------------------------------------------------------------------------
%
\begin{frame}{Tail call optimisation}

  The same, with a garbage collector that collects after each
  replacement:
  \begin{center}
    \includegraphics[bb=71 644 378 721]{sum375_pop1}
  \end{center}

\end{frame}

% ------------------------------------------------------------------------
%
\begin{frame}{Tail call optimisation}

  Let us consider an alternate version of `\fun{sum}':
  \begin{equation*}
    \begin{array}{@{}r@{\;}l@{\;}l@{}}
      \fun{sum}_1(\cons{n}{t}) & \xrightarrow{\gamma}
      & \fun{sum}_0(n,t).\\
      \fun{sum}_0(n,\el) & \xrightarrow{\delta} & n;\\
      \fun{sum}_0(n,\cons{m}{s}) & \xrightarrow{\epsilon}
      & \fun{sum}_0(n+m,s).
    \end{array}
  \end{equation*}

  For \(\fun{sum}_1([3;7;5])\) we have the evaluation
  \begin{center}
    \includegraphics[bb=71 641 342 722]{sum0_375_0}
  \end{center}

\end{frame}

% ------------------------------------------------------------------------
%
\begin{frame}{Tail call optimisation}

  \begin{center}
    \includegraphics[bb=71 641 346 721]{sum0_375_1}
    \includegraphics[bb=71 639 384 721]{sum0_375_2}
  \end{center}

\end{frame}

% ------------------------------------------------------------------------
%
\begin{frame}{Tail call optimisation}

  Notice that, in the case of \fun{sum}$_0$, we found the result (15)
  at the end of the first phase (of growing the call stack).

  \bigskip

  Therefore, there is no need for a second phase, which replaces the
  pending calls by their values: we can discard all those calls at
  once after the first phase.

  \bigskip

  Or, to see it even more clearly, let us rewrite the same first phase
  and assume that the collector reclaims immediately the useless
  pending calls:
  \begin{center}
    \includegraphics[bb=71 640 326 721]{sum0_375tco0}
  \end{center}

\end{frame}

% ------------------------------------------------------------------------
%
\begin{frame}{Tail call optimisation}

  \begin{center}
    \includegraphics[bb=71 628 350 721]{sum0_375tco1}
    \includegraphics[bb=71 639 366 721]{sum0_375tco2}
  \end{center}

\end{frame}

% ------------------------------------------------------------------------
%
\begin{frame}{Tail call optimisation}

  As we see now, we can keep the size of the call stack constant, if
  we reclaim the top at each new rewrite.

  \bigskip

  This optimisation is called \textbf{tail call optimisation}.

  \bigskip

  Intuitively, a call in \textbf{tail position} is a call whose value
  is the value of the previously rewritten call.

  \bigskip

  For example, the call \(f(x)\) in \(1 + f(x)\) is not in tail
  position, because its value needs to be incremented.

  \bigskip

  Most compilers of functional languages detect and implement that
  optimisation that saves memory, including the \Clang compiler of
  \textsf{GCC}.

\end{frame}
